\documentclass[11pt,a4paper]{article}
\usepackage[margin=25mm]{geometry}
\usepackage{amsmath,amssymb,amsthm,booktabs,hyperref}
\newtheorem{theorem}{Theorem}\newtheorem{lemma}[theorem]{Lemma}
\newcommand{\Aut}{\operatorname{Aut}}\newcommand{\Irr}{\operatorname{Irr}}
\title{A character bound for simultaneous inversion}
\author{Human Compression project}\date{}
\begin{document}\maketitle
\begin{abstract}
A generating pair with involutory product has at most one simultaneous
inverter. We express such pairs as triples of involutions and bound their
number using two class multiplication counts. Together with the previously
proved generating-pair bounds, this proves the existence of chiral maps
for 100 of the 128 types asserted in Lemma 5.1 of the fixed O'Nan source.
The proof uses the named character tables as explicit premises and does
not enumerate generating pairs in the O'Nan group.
\end{abstract}

\section{The inversion correspondence}
Let $G$ be a finite centreless group and put $A=\Aut(G)$. Identify $G$
with its group of inner automorphisms, using $g:x\mapsto gxg^{-1}$.
Products below are products in $A$, and $bRb^{-1}=b(R)$ under this
identification. For fixed $p,q>2$, let
\[
 E_{p,q}=\{(R,S)\in G^2: |R|=p,\ |S|=q,\ |RS|=2,
                    \ \langle R,S\rangle=G\}.
\]
The action of $A$ on this set is free: an automorphism fixing both
generators fixes every element of $G$. Call a pair reflexible if some
$b\in A$ inverts both $R$ and $S$, and chiral otherwise.

\begin{lemma}\label{triple}
A reflexible pair has a unique simultaneous inverter $b$, and $b$ has
order two. The assignments
\[
 (R,S)\longmapsto (a,b,c)=(Rb,b,bS),\qquad
 (a,b,c)\longmapsto (ab,bc)
\]
give mutually inverse $A$-equivariant correspondences between reflexible
pairs and triples of nonidentity involutions satisfying
\[
 ab,bc\in G,\quad |ab|=p,\quad |bc|=q,\quad |ac|=2,
                    \quad\langle ab,bc\rangle=G.
\]
In these triples $a$ and $c$ commute and are distinct.
\end{lemma}
\begin{proof}
If $b$ and $b'$ both invert the generators, then $(b')^{-1}b$ fixes
both generators, so is the identity. Also $b^2$ fixes both generators,
and $b$ cannot be the identity since $p,q>2$. Thus $b^2=1$.
The identities $bRb=R^{-1}$ and $bSb=S^{-1}$ give
$(Rb)^2=(bS)^2=1$. Neither product is the identity, since otherwise
$R$ or $S$ would have order two. We have $ab=R$, $bc=S$ and
$ac=RS$. Two involutions have product of order two precisely when they
are distinct and commute.

Conversely, from such a triple put $R=ab$ and $S=bc$.
Conjugation by $b$ sends $ab$ to $ba=(ab)^{-1}$ and $bc$ to
$cb=(bc)^{-1}$. Hence $b$ is the unique inverter of this generating
pair. The two displayed assignments are inverse. They use only products
and the unique inverter, so commute with the action of $A$.
\end{proof}

Assume now that $[A:G]=2$ and that each of the two cosets of $G$
contains exactly one $A$-conjugacy class of involutions. Denote these
classes by $C_0\subset G$ and $C_1\subset A\setminus G$.
In Lemma~\ref{triple}, $ab,bc\in G$ forces all three involutions
to lie in the same coset and therefore in the same class $C_i$.
For $a_i\in C_i$ define
\[
 \kappa_t(i)=|\{b\in C_i:|a_i b|=t\}|,
 \qquad h_i=|C_A(a_i)|.
\]
In particular, $\kappa_2(i)$ counts the involutions in $C_i$
which commute with $a_i$, excluding $a_i$ itself. It does not count
involutions in the other coset.

\begin{theorem}\label{bound}
Let $r_i(p,q)$ be the number of $A$-orbits of generating pairs of type
$(p,q)$ whose inverter belongs to $C_i$. Then
\[
 r_i(p,q)\leq
 U_i(p,q):=\left\lfloor
 \frac{\kappa_2(i)\min\{\kappa_p(i),\kappa_q(i)\}}{h_i}
 \right\rfloor.                                      \tag{1}
\]
Consequently, if $L(p,q)\leq |E_{p,q}|/|A|$, the number of chiral
orbits is at least $\max\{0,L(p,q)-U_0(p,q)-U_1(p,q)\}$.
\end{theorem}
\begin{proof}
Choose $a\in C_i$, then $b\in C_i$ with $|ab|=p$, and then
$c\in C_i$ with $|ac|=2$. This gives
$|C_i|\kappa_p(i)\kappa_2(i)$ triples before imposing $|bc|=q$
or generation. It is therefore an upper bound for the triples in
Lemma~\ref{triple}. Reversing $(a,b,c)$ to $(c,b,a)$ exchanges $p$
and $q$, giving the same bound with $\kappa_q(i)$. The action on
generating triples is free by the lemma and the free action on pairs.
Divide the smaller bound by $|A|=|C_i|h_i$, then take its integer
part. The two classes of inverters are disjoint and exhaustive, which
proves the final assertion.
\end{proof}
These are counts for ordered generators with prescribed orders. We
display one orientation $p\leq q$, as in the retained count components.
We do not divide by reversal or duality, including when $p=q$.
The lemma makes no assertion about the intersection property of the
three involutions; a string C-group census has that additional condition.

\section{Obtaining the bound from characters}
Write $D_t=\{x\in G:|x|=t\}$, viewed as an $A$-invariant subset
of $A$. It is closed under inverses. For $\chi\in\Irr(A)$ put
\[
 v_{\chi,t}=\sum_{x\in D_t}\chi(x),\qquad
 w_{\chi,i}=|C_i|\chi(a_i).
\]
\begin{lemma}\label{characters}
With the above conventions,
\[
 \kappa_t(i)=\frac1{|A|}\sum_{\chi\in\Irr(A)}
       \frac{\chi(a_i)w_{\chi,i}v_{\chi,t}}{\chi(1)}.       \tag{2}
\]
\end{lemma}
\begin{proof}
In the complex group algebra, the coefficient of $a_i$ in
$\bigl(\sum_{b\in C_i}b\bigr)\bigl(\sum_{x\in D_t}x\bigr)$
counts the $b\in C_i$ for which $b^{-1}a_i\in D_t$.
Since $b$ is an involution and $|ba_i|=|a_i b|$, this is
$\kappa_t(i)$. On an irreducible module the two central sums act
by $w_{\chi,i}/\chi(1)$ and $v_{\chi,t}/\chi(1)$.
The coefficient of $a_i$ in the corresponding central idempotent is
$\chi(1)\chi(a_i^{-1})/|A|$. As $a_i^{-1}=a_i$, their product
gives (2). This is the usual class multiplication formula, specialized
to one involution class and an order set.
\end{proof}

For $G=O'N$ we use the retained named tables \texttt{ON} and
\texttt{ON.2}, their stored fusion, and the identification
$A=\Aut(O'N)$. These table identifications, class coverage and fusion
are borrowed premises. The acquired data say that the socle is the
union of classes 1 through 25 of \texttt{ON.2}, and its two involution
classes are 2 and 26. Their class sizes and centralizers are
\[
 (|C_0|,h_0)=(2857239,322560),\qquad
 (|C_1|,h_1)=(2624832,351120).
\]
In particular the outer denominator is the full centralizer in $A$.
The values in (2) are computed with exact rational arithmetic from the
retained vectors $(\chi(1),\chi(a_i),w_{\chi,i},v_{\chi,t})$.
The reader checks integrality, column norms, the character degree sum,
the class-size equations and the complete partition
$\sum_{t>1}\kappa_t(i)=|C_i|-1$. The identity product occurs exactly
once, when $b=a_i$. These consistency checks do not replace the named
character-table premises.
\begin{center}\begin{tabular}{rrr}\toprule
$t$&$\kappa_t(0)$&$\kappa_t(1)$\\\midrule
2&1750&1463\\
3&20608&5852\\
4&42560&46816\\
5&44800&58520\\
6&120960&29260\\
7&46080&200640\\
8&241920&0\\
10&80640&175560\\
11&161280&175560\\
12&161280&58520\\
14&161280&175560\\
15&322560&117040\\
16&645120&0\\
19&483840&526680\\
20&0&351120\\
28&322560&351120\\
31&0&351120\\
\bottomrule\end{tabular}\end{center}


\section{The recovered part of the source assertion}
We insert in Theorem~\ref{bound} the previously accepted lower bounds
$L(p,q)$ from maximal-subgroup subtraction, strengthened by the exact
columns for orders 11, 19 and 31. These are independent of the printed
Table 1. The resulting lower bound is positive for exactly the following
100 types among the source's 128 asserted chiral types.
\begin{center}\begin{tabular}{rl}\toprule
$p$&values of $q$\\\midrule
3&20, 31\\
4&8, 11, 16, 19, 20, 31\\
5&8, 11, 16, 19, 20, 28, 31\\
6&8, 10, 11, 15, 16, 19, 20, 28, 31\\
7&8, 10, 11, 12, 14, 15, 16, 19, 20, 28, 31\\
8&8, 10, 11, 12, 14, 15, 16, 19, 20, 28, 31\\
10&10, 11, 12, 14, 15, 16, 19, 20, 28, 31\\
11&11, 12, 14, 15, 16, 19, 20, 28, 31\\
12&14, 15, 16, 19, 20, 28, 31\\
14&14, 15, 16, 19, 20, 28, 31\\
15&15, 16, 19, 20, 28, 31\\
16&16, 19, 20, 28, 31\\
19&19, 20, 28, 31\\
20&20, 28, 31\\
28&28, 31\\
31&31\\
\bottomrule\end{tabular}\end{center}

For example, type $(3,20)$ has generating lower bound
$44$, inner reflexible upper bound $0$ and outer reflexible upper bound $24$; hence at least $20$ of its orbits are chiral.
Both involution classes are included in this subtraction.
Summing the 136 lower bounds gives 690899 chiral generating-pair orbits
under these conventions. This is a lower bound, not an exact census.

The complete numerical ledger records all 136 applications of (1).
Of the positive types, the following 28 asserted chiral types still
have zero lower bound from this method:
\begin{center}\begin{tabular}{rl}\toprule
$p$&values of $q$\\\midrule
3&8, 10, 11, 14, 16, 19, 28\\
4&5, 6, 7, 10, 12, 14, 15, 28\\
5&5, 6, 7, 10, 12, 14, 15\\
6&6, 7, 12, 14\\
7&7\\
12&12\\
\bottomrule\end{tabular}\end{center}

The exceptional type $(3,15)$ also remains undecided by this argument;
we do not recover the source's claim that all six of its pairs are
reflexible, with four inner and two outer inverters. No new exact
entry of Table 1 is obtained here: its accepted status remains
55 of 136, with 81 unresolved. The positive inner and outer spectra
and the remaining assertions of Section 6 require separate work.

The involution correspondence and bound (1) are the present argument.
The class multiplication formula, the named character tables and the
earlier generating lower bounds are reused inputs. Lemma 5.1 of
Connor and Leemans~\cite{CL} is the source assertion being partially
reproved; its search results are not premises of the bound. A literal
check in $A_5\lhd S_5$ verifies the correspondence, unique inverters,
free orbit sizes and the bound for every occurring type, including
equal orders. This finite check supplements the general proof and
does not certify the O'Nan table identifications. No new open
conjecture or whole-original completion is claimed.
\begin{thebibliography}{9}
\bibitem{CL} T. Connor and D. Leemans, \emph{Algorithmic enumeration of
regular maps}, Ars Math. Contemp. 10 (2016), 211--222, Section 5.
\bibitem{CT} T. Breuer, \emph{The GAP Character Table Library},
retained installed edition and tables \texttt{ON}, \texttt{ON.2};
exact library bytes and fusion export bound in the acquisition receipt.
\end{thebibliography}
\end{document}
